\session{11}{11/7}{2限}{3}{データに基づく意思決定}

\goals{
  \item データに基づく意思決定（データドリブン）の意味と、勘や経験との関係を説明できる。
  \item 人の判断に入り込む認知バイアスの代表例を挙げ、対策を考えられる。
  \item AIの推奨と向き合う設計（人が最終判断に関与する設計）と、意思決定の説明責任を理解する。
}

\guidesec{解説}

\topic{データに基づく意思決定}
データドリブンな意思決定とは、勘や経験だけでなく、データの分析結果を根拠として判断することです。
勘や経験を否定するものではありません。経験は仮説を生み、データは仮説を検証します。
データだけで決めようとすると、データにない要因（従業員の士気、顧客との関係、倫理）を見落とし、
データの偏りや分析の前提（第10回）を疑わなくなる「データ万能主義」に陥ります。

データに基づく意思決定の流れは、第2回のサイモンの3段階に対応します。情報活動でデータを集めて問題を定義し、
設計活動で代替案とその予測される結果をデータで比べ、選択活動で評価基準に照らして選びます。
そして実行後に結果を測り、次の意思決定に活かします（再検討）。この循環を回すには、
意思決定の時点で「何を期待していたか」を記録しておくことが必要です。

\topic{人の判断のバイアス}
人の判断には、系統的な偏り（認知バイアス）があります。代表的なものを挙げます。
\par\noindent\begin{gtabh}{colspec={Q[l,wd=30mm] X[l] X[l]}}
  バイアス & 内容 & 経営での例 \\
  確証バイアス & 自分の考えを支持する情報ばかりを集め、反する情報を軽視する & 新規事業に賛成の資料だけを集める \\
  アンカリング & 最初に示された数値や情報に判断が引きずられる & 最初の見積額が交渉の基準になる \\
  利用可能性ヒューリスティック & 思い出しやすい事例（最近の出来事、印象的な事件）を過大評価する & 直近のクレーム1件で方針を変える \\
  サンクコスト & すでに投じた費用に引きずられ、撤退できない & 効果の出ないシステムに追加投資する \\
  過信 & 自分の判断や予測の正確さを過大評価する & 予測の幅を狭く見積もる \\
  集団思考 & 集団の和を優先し、異論が出なくなる & 会議で社長の案に誰も反対しない \\
  自動化バイアス & 機械やAIの出力を、人の判断より無条件に信頼する & AIの推奨をそのまま採用する \\
\end{gtabh}
対策の基本は、判断の手順を決めておくこと（評価基準を先に決める、反対意見を必ず出す役を置く）、
判断を記録して後から検証すること、そして自分のバイアスを知っておくことです。生成AIに「この判断に反する証拠は」
「反対の立場ならどう言うか」と尋ねることは、確証バイアスと集団思考への実用的な対策になります。

\topic{AIの推奨とどう向き合うか}
AIがデータから推奨（この顧客に提案せよ、この発注量にせよ）を出す場面が増えています。
AIの推奨をどう扱うかは、設計の問題です。
\begin{itemize}
  \item \textbf{Human-in-the-loop}　AIの処理の途中や最終段階に人の確認・判断を組み込む。人が承認しなければ実行されない。
  \item \textbf{Human-on-the-loop}　AIが自動で実行するが、人が監視し、必要なら介入・停止できる。
  \item \textbf{Human-out-of-the-loop}　人が関与せず、AIが自動で判断・実行する。
\end{itemize}
影響が大きく、取り消しが難しく、個人の権利に関わる判断ほど、人の関与を厚くします（採用、与信、医療など）。
量が多く、影響が小さく、取り消しが容易な判断は、自動化して人は監視に回ります（商品の推薦など）。
人が関与する設計にしても、自動化バイアスによって形だけの確認になれば意味がありません。
人が実質的に判断できるように、AIの推奨の根拠を示す、人が反対しやすい手順にする、人の判断とAIの推奨の一致率を監視する、といった工夫が必要です。

\topic{意思決定の説明責任}
AIの推奨を採用して判断した結果の責任は、判断した人・組織にあります。AIは法的な責任の主体ではありません。
説明責任を果たすには、意思決定の時点で、何を根拠に（どのデータ、どの分析、どのAIの推奨）、何を期待して、
誰が決めたかを記録しておく必要があります。1枚の\textbf{意思決定メモ}にまとめておくと、後からの検証と、
関係者への説明が容易になります。

\topic{キーワード}
\par\noindent\begin{keywords}{}
  データドリブン & 勘や経験に加えて、データの分析結果を根拠に判断すること。 \\
  認知バイアス & 人の判断に系統的に入り込む偏り。確証バイアス、アンカリングなど。 \\
  自動化バイアス & 機械やAIの出力を人の判断より無条件に信頼してしまう傾向。 \\
  Human-in-the-loop & AIの判断過程に人が関与し、最終判断や確認を人が行う設計。 \\
  説明責任 & 判断の根拠と経緯を関係者に説明できる状態にしておく責任。 \\
  意思決定メモ & 問題、根拠、代替案、決定、期待する結果、責任者を1枚にまとめた記録。 \\
\end{keywords}

\guidesec{演習課題}

\exercise{1}{AI演習：分析結果をもとに、AIと共同で1枚の意思決定メモを作成する}
\instr{第9・10回の商品別売上の分析結果をもとに、「商品Bについて来期どうするか」の意思決定メモを、次のプロンプトでAIと共同で作成してください。AIの案に対して、必ず「この案に反する証拠は」「反対の立場なら」と尋ね、その結果をメモに反映してください。メモの最後に、AIの寄与と自分の判断を分けて書きます。}
\begin{promptbox}[意思決定メモの作成]
あなたは経営企画の担当者です。次の分析結果をもとに、「商品Bについて来期どうするか」の意思決定メモを作ります。メモの項目は、(1) 問題の定義、(2) 根拠となるデータと分析、(3) 代替案（3つ以上）とそれぞれの予想される結果、(4) 評価基準、(5) 推奨案と理由、(6) 想定されるリスクと反対意見、(7) 期待する結果と、それを測る指標、(8) 決定者と決定日、です。まず (1)〜(4) を書いて、私の意見を聞いてから (5) 以降に進んでください。
（分析結果：第9回の表と、第10回の予測）
\end{promptbox}

\exercise{2}{バイアスを見つける}
\instr{次の会議の記録から、認知バイアスが疑われる発言を3つ抜き出し、どのバイアスか、どう対策するかを書いてください。}
\begin{point}{会議の記録（抜粋）}
社長「商品Bは私が立ち上げた商品だ。ここまで投資してきたのに、今やめるわけにはいかない」。
営業部長「先週、大口の顧客から商品Bを褒められました。まだ需要はあります」。
企画担当「AIの予測では来期も減少です。AIが言うのだから間違いないでしょう」。
経理部長「最初の計画では商品Bは年間400の売上を見込んでいました。それに比べれば今の290は惨憺たるものです」。
全員「社長がそう言うなら、継続で」。
\end{point}

\exercise{3}{人の関与を設計する}
\instr{次の3つのAIの利用場面について、Human-in-the-loop／on-the-loop／out-of-the-loop のどれが適切か、影響の大きさ・取り消しの容易さ・件数の観点から判定し、人が実質的に判断できるようにする工夫を1つずつ書いてください。}
\begin{itemize}
  \item (A) EC サイトでの商品の推薦
  \item (B) 採用選考での書類の一次評価
  \item (C) 店舗の日配品の自動発注
\end{itemize}

\guidesec{模範解答}

\begin{answer}{演習1}
意思決定メモの例（要約）：
\par\noindent\begin{gtab}{colspec={Q[l,wd=30mm] X[l]}}
  問題の定義 & 商品Bの売上が2023年の390から2026年の290へ減少（約26\%減）。合計は伸びているが、商品Bの減少が続けば2028年には商品Cと合わせた成長が止まる可能性。 \\
  根拠 & 第9回の売上表、第10回の線形予測（2027年 285）。限界：年次5点のみ、価格・競合・顧客別の情報なし。 \\
  代替案 & (a) 現状維持、(b) 価格と販促の見直しで回復を狙う、(c) 商品Cへ経営資源を移し、商品Bは縮小、(d) 商品Bの顧客を調査してから決める（決定を3か月延期）。 \\
  評価基準 & 来期の合計売上への影響、必要な投資、実行の確実さ、顧客への影響、判断の取り消しやすさ。 \\
  推奨案と理由 & (d) を採用し、その間 (b) の小規模な試行を行う。理由：減少の原因（商品Cへの移行か、商品B固有の問題か）が分からない段階で (c) を選ぶと取り消しが難しい。 \\
  リスクと反対意見 & 延期の間に減少が加速する。反対意見：「原因調査は時間の無駄で、商品Cへの集中が合理的」。これに対しては、試行の期限を3か月と区切る。 \\
  期待する結果と指標 & 3か月後に減少の原因を特定。試行した販促の対象店舗での商品Bの前年同月比。 \\
  決定者と決定日 & 経営会議（議長：社長）、2026年11月7日。 \\
\end{gtab}
AIの寄与：項目の構成、代替案 (a)〜(c) の列挙、リスクの列挙。自分の判断：代替案 (d) の追加（AIは「調査してから決める」案を出さなかった）、
評価基準に「取り消しやすさ」を加えたこと、推奨案の選択。検証：予測値は第10回の演習で自分で再計算したものを使った。
\end{answer}

\begin{answer}{演習2}
\begin{itemize}
  \item 社長「ここまで投資してきたのに」：サンクコスト。対策：これまでの投資は判断から切り離し、「今後の投資と期待される結果」だけで評価する基準を先に決める。
  \item 営業部長「先週、大口の顧客から褒められた」：利用可能性ヒューリスティック（印象的な1件の過大評価）。対策：顧客全体のデータ（購買の継続率、解約率）で確かめる。
  \item 企画担当「AIが言うのだから間違いない」：自動化バイアス。対策：予測の前提と限界（第10回）を確認し、AIの推奨を「意見の1つ」として扱う。
  \item 経理部長「最初の計画の400に比べれば」：アンカリング（当初の計画値が基準になっている）。対策：市場の変化を踏まえた現実的な基準で評価し直す。
  \item 全員「社長がそう言うなら」：集団思考。対策：決定の前に反対意見を出す役を決める。匿名で意見を集める。
\end{itemize}
（5つ挙げたが、課題では3つで十分です。）
\end{answer}

\begin{answer}{演習3}
\begin{itemize}
  \item (A) out-of-the-loop（人は監視）：影響が小さく、件数が膨大で、推薦が外れても顧客は無視できる。工夫：推薦の成果（クリック率、購入率）と苦情を人が定期的に監視し、不適切な推薦（差別的な組み合わせなど）を検知する基準を決める。
  \item (B) in-the-loop：個人の権利に関わり、影響が大きく、取り消しが難しい。工夫：AIの評価を「順位」ではなく「確認すべき点の指摘」として示し、人が独立に評価してから AI の評価と比べる。AIの評価と人の判断の一致率を監視し、形だけの確認になっていないか確かめる。学習データの偏り（第13回）を定期的に点検する。
  \item (C) on-the-loop：件数が多く、取り消し（追加発注、値引き）が比較的容易だが、影響は廃棄と欠品という形で出る。工夫：発注量が通常の範囲を外れたときだけ人に通知し、承認を求める。天候やイベントなどAIが知らない情報を担当者が追加できる手順を用意する。
\end{itemize}
\end{answer}

\guidesec{出席確認クイズの解説}
講義サイトの出席確認ページ（第11回）の5問です。
% 出席確認クイズ 第11回　データに基づく意思決定
%   attendance/attendance11.html から生成（解説は本ガイド用に加筆）。

\quizq{1}{「確証バイアス」の説明として適切なものはどれでしょうか?}%
  {損失を過大に嫌う傾向}%
  {\correct{自分の考えを支持する情報ばかりを集めてしまう傾向}}%
  {多数派の意見に合わせる傾向}%
  {最初に見た数値に判断が引きずられる傾向}%
  {正解は (b)。確証バイアスは、自分の仮説に都合のよい情報を重視し、反する情報を軽視してしまう傾向です。(a) は損失回避、(c) は同調（集団への追随）、(d) はアンカリングの説明です。AIに「この考えに反する証拠は」と尋ねるのは、確証バイアスへの対策になります。}

\quizq{2}{「アンカリング」の説明として適切なものはどれでしょうか?}%
  {\correct{最初に提示された数値や情報に判断が引きずられる}}%
  {自分の考えを支持する情報を集める}%
  {直近の出来事を過大評価する}%
  {集団で極端な結論に至る}%
  {正解は (a)。アンカリングは、最初に示された数値や情報（錨）を基準にして、その後の判断が偏ってしまう現象です。(b) は確証バイアス、(c) は利用可能性ヒューリスティック、(d) は集団極性化の説明です。AIが最初に出した数値にも引きずられやすいので注意が必要です。}

\quizq{3}{「Human-in-the-loop」の説明として適切なものはどれでしょうか?}%
  {\correct{AIの判断過程に人が関与し、最終判断や確認を人が行う設計}}%
  {人がまったく関与しない全自動化}%
  {AIを使わないこと}%
  {人がAIの学習データになること}%
  {正解は (a)。Human-in-the-loop は、AIの処理の途中や最終段階に人の確認・判断を組み込む設計の考え方です。(b) は人が関与しない全自動化（human-out-of-the-loop）です。影響の大きい判断ほど人の関与を厚くします。}

\quizq{4}{「データに基づく意思決定(データドリブン)」の説明として適切なものはどれでしょうか?}%
  {データがあれば人の判断は不要}%
  {\correct{勘や経験だけでなく、データの分析結果を根拠に判断する}}%
  {データが多いほど無条件に正しい}%
  {データを無視して直感で決める}%
  {正解は (b)。データドリブンな意思決定は、勘や経験に加えてデータの分析結果を根拠として判断することです。(a)(c) はデータ万能主義で、データの偏りや分析の前提を無視しています。(d) はデータドリブンの逆です。}

\quizq{5}{AIの推奨に従って判断した場合、その結果の責任は誰にあるでしょうか?}%
  {\correct{最終的な判断をした人・組織}}%
  {誰にもない}%
  {AIそのもの}%
  {AIの開発者だけ}%
  {正解は (a)。AIは判断を支援する道具であり、AIの推奨を採用した最終的な責任は判断した人・組織にあります。(c) AI は法的な責任主体ではなく、(d) 開発者の責任は製品の欠陥など限られた場面に限られます。説明責任を果たすために、判断の根拠と経緯を記録しておくことが重要です。}

